\section{Production plan}\label{sec:production}

\subsection{Folder layout}
\begin{verbatim}
packages_book/
  README.md              what this folder is, how to build
  plan/                  this document (main, macros, sec*.tex, refs.bib)
  book/                  created when writing starts
    main.tex  macros.tex  refs.bib
    chapters/ch01.tex ... ch14.tex  appA.tex
    figures/  tables/  listings/    generated, never edited by hand
    scripts/                        one script per figure or table
\end{verbatim}
The book reuses the conventions of \code{paper/}: a shared macro file, numbers generated into macros, figures built by scripts, and \texttt{latexmk} for the build.

\subsection{Order of work}
\begin{enumerate}
\item \textbf{Set-up.} Book skeleton, macro file, numbers pipeline, bibliography, the nine items of \Cref{tab:todo}.
\item \textbf{Part II core first} (Chapters~\ref{ch:elem}, \ref{ch:nonlin}, \ref{ch:batch}, \ref{ch:contact}), because it holds the most measured material and the contact chapter is the book's distinctive contribution.
\item \textbf{Part I and Chapter~\ref{ch:dyn}}, written after the code it describes is fixed.
\item \textbf{Part III} in the order of \Cref{fig:deps}.
\item \textbf{Part IV and the appendix}, then index, notation list and the anchor-to-path table.
\end{enumerate}
Each chapter is reviewed against the repository at its release tag, and every listing is executed in the build.

\subsection{Risks}
\begin{itemize}
\item \textbf{Code drift.} The packages still change. Mitigation: pin a tag per edition and cite by importable name.
\item \textbf{Unmeasured claims.} The new gallery scripts (\code{fea_09}, \code{fea_10}, \code{rom_07}) were run only in the development sandbox; they need a run on the author's machine before their numbers enter the book.
\item \textbf{Scope creep.} The package has many modules that the plan leaves out (topology optimisation, adaptivity, iterative solvers, harmonic-balance hyper-reduction, friction models). They appear only as pointers.
\item \textbf{Length.} 440 pages is a ceiling; Chapters~\ref{ch:elem} and~\ref{ch:nonlin} are the first to shorten.
\end{itemize}

\subsection{Decisions taken}
The five open questions of draft~1 were answered by the author.
\begin{enumerate}
\item \textbf{Companion to a pinned release.} The book is a companion to a tagged release of the repository; it is not a stand-alone text. Every listing, number and code anchor refers to that release. The current versions are \fea{} 1.0.1 and \rom{} 0.2.0; the tag itself is to be created by the author when the text is frozen, and all results quoted in the book are re-run at that tag.
\item \textbf{FD-PINN reproduction stays in Chapter~\ref{ch:verif}} as a verification case study, presented as a check of the reproduction procedure and not as a neural-network method.
\item \textbf{One paragraph on the neural-operator code} in Chapter~\ref{ch:nonint}; no neural-network chapter.
\item \textbf{Format.} A4, class \texttt{book}, in the layout of \code{paper/macros.tex}.
\item \textbf{Appendix~A (differentiable assembly) stays.}
\end{enumerate}

\subsection{Blueprint review: items adopted}\label{sec:blueprint}
An outside blueprint for a longer, from-scratch textbook (24 chapters, Parts on FEM theory, MOR theory and integrated workflows) was compared with this plan. Its Parts on linear algebra, continuum mechanics, weak forms and element derivations are out of scope for a companion book (decision~1 above), as are its 24-chapter structure, the 12-month schedule, a printed API reference and extended GPU material. The following items are adopted. They are applied in one editing pass after Chapters~\ref{ch:nonint} and~\ref{ch:hyper} are written, so that no chapter is revised twice; status is updated here.
\begin{enumerate}
\item \textbf{Learning objectives.} Three to four ``after this chapter the reader can \dots'' items added to the chapter box of every chapter. \emph{Status: not done.}
\item \textbf{Three exercise types.} Each chapter's exercises contain at least one derivation, one implementation and one critical-interpretation exercise, tagged as such. \emph{Status: not done.}
\item \textbf{``How to read the verification sections''.} A short front-matter page defining unit test, verification and validation (no result in this book is validation against experiments; a passing test is not proof of general correctness) and listing what each reported number states: model and loading, solver settings, reference value, error definition, release and machine. A back-matter appendix gives the convergence-study template and the tolerances used. \emph{Status: not done.}
\item \textbf{Fuller ROM reporting} in Chapters~\ref{ch:pod}, \ref{ch:freq} and~\ref{ch:verif}: offline, online and reconstruction time reported separately; each speed-up states whether it includes preprocessing; a local quantity (a reaction or peak stress) reported next to the global relative error. \emph{Status: partly in Chapter~\ref{ch:pod} (offline and online separated, dense baseline stated); local quantities and reconstruction time not done.}
\item \textbf{Benchmark discrepancy.} The README and user guide give 221.8$\times$ (1653.23~ms against 7.45~ms, 500-point sweep, maximum error $1.51\times10^{-3}$); a figure of 254$\times$ from an earlier manuscript does not occur anywhere in the repository. Chapter~\ref{ch:verif} quotes 221.8$\times$ against the tagged release and gives 254$\times$ only with its own conditions; the two are never merged. \emph{Status: open; needs the source of 254$\times$ and a release tag.}
\item \textbf{One complete integrated run} (full-order model, snapshots, reduced model, held-out check, reconstruction) in Chapter~\ref{ch:verif}. \emph{Status: not done (chapter not written).}
\item \textbf{Maturity labels.} A stable / experimental / not-covered column in each chapter's code-map table. \emph{Status: not done.}
\item \textbf{Glossary and notation appendix}, and a table mapping each chapter to its script and to the package tests that cover it. \emph{Status: not done.}
\end{enumerate}
\subsection{Progress}
The \code{book/} skeleton exists (A4 \texttt{book} class, macro file, bibliography, stubs for all chapters and Appendix~A, build with \texttt{latexmk}). Chapters~\ref{ch:arch}, \ref{ch:mesh}, \ref{ch:elem}, \ref{ch:dyn}, \ref{ch:nonlin}, \ref{ch:batch}, \ref{ch:contact}, \ref{ch:pod}, \ref{ch:bt}, \ref{ch:freq}, \ref{ch:cms} and~\ref{ch:nonint} are written in full, and Chapter~\ref{ch:contact} is the template: it has the chapter box, formulation, algorithm box, tested listings, figures and tables generated by \code{book/scripts/ch07_assets.py}, a limits box, a failure-mode table, a summary and exercises. It currently has 11 pages against a budget of 36; the budget is an upper bound, and the plan's later chapters should be sized by their content. Each has its own assets script (\code{ch01_assets.py}, \code{ch02_assets.py}, \code{ch03_assets.py}, \code{ch04_assets.py}, \code{ch05_assets.py}, \code{ch06_assets.py}, \code{ch07_assets.py}, \code{ch08_assets.py}, \code{ch09_assets.py}, \code{ch10_assets.py}, \code{ch11_assets.py}, \code{ch12_assets.py}), so Parts~I and~II are complete and Part~III is half written. The next chapters to write are Chapter~13, then Chapter~14 and Appendix~A.
